% ------------------------------------------------- worked example, one polynomial
\begin{table}[!tbh]
\centering\scriptsize\setlength{\tabcolsep}{4pt}
\begin{tabular}{|l|p{.50\linewidth}|l|}
\hline
\textbf{Code} & \textbf{Exact action on $P(x) = 2 + 5x + 3x^2 + 4x^3$} & \textbf{Result} \\
\hline
SH1 & shift of the argument, $P(x) \to P(x+1)$: $2 + 5(x+1) + 3(x+1)^2 + 4(x+1)^3 = 14 + 23x + 15x^2 + 4x^3$ & $[3, 1, 4, 4]$ \\ \hline
SC2 & scale of the argument, $P(x) \to P(2x)$: the coefficient at $x^i$ is multiplied by $2^i$, giving $[2, 10, 12, 32]$ & $[2, 10, 1, 10]$ \\ \hline
REV & reversal of the coefficient vector. This is neither a reversal of digits nor a substitution of $-x$ for $x$: $[c_0, c_1, c_2, c_3] \to [c_3, c_2, c_1, c_0]$ & $[4, 3, 5, 2]$ \\ \hline
AC1 & one is added to the free coefficient only: $[2+1, 5, 3, 4]$ & $[3, 5, 3, 4]$ \\ \hline
AX1 & one is added to the coefficient at $x$ only: $[2, 5+1, 3, 4]$ & $[2, 6, 3, 4]$ \\ \hline
\end{tabular}
\caption{All five transformations on one starting polynomial over $F_{11}$, with the
state $[2, 5, 3, 4]$. Coefficients are reduced modulo 11 after every operation. AC1
and AX1 look similar and change different coordinates, so after REV or
SH1 they lead to different final states. Every row was re-executed with the same
deterministic interpreter that the experiments use.}
\label{tab:opsexample}
\end{table}
